\documentclass[11pt]{article}
\usepackage[T1]{fontenc}
\usepackage[utf8]{inputenc}
\usepackage{lmodern}
\usepackage[margin=1.05in]{geometry}
\usepackage{amsmath,amssymb,amsthm,mathtools}
\usepackage{microtype,booktabs,array}
\usepackage{xcolor}
\usepackage{listings}
\usepackage[colorlinks=true,linkcolor=blue!45!black,citecolor=blue!45!black,urlcolor=blue!45!black]{hyperref}
\usepackage[nameinlink,capitalise,noabbrev]{cleveref}
\numberwithin{equation}{section}
\newtheorem{theorem}{Theorem}[section]
\newtheorem{proposition}[theorem]{Proposition}
\newtheorem{lemma}[theorem]{Lemma}
\newtheorem{corollary}[theorem]{Corollary}
\newtheorem{conjecture}[theorem]{Conjecture}
\theoremstyle{definition}
\newtheorem{definition}[theorem]{Definition}
\newtheorem{construction}[theorem]{Construction}
\theoremstyle{remark}
\newtheorem{remark}[theorem]{Remark}
\newtheorem{example}[theorem]{Example}
\newcommand{\N}{\mathbb N}
\crefname{construction}{Construction}{Constructions}
\crefname{lemma}{Lemma}{Lemmas}
\crefname{proposition}{Proposition}{Propositions}
\crefname{conjecture}{Conjecture}{Conjectures}
\crefname{example}{Example}{Examples}
\crefname{remark}{Remark}{Remarks}
\crefname{definition}{Definition}{Definitions}
\crefname{corollary}{Corollary}{Corollaries}

\lstdefinelanguage{lean}{
  morekeywords={def,theorem,lemma,by,fun,let,in,if,then,else,match,with,structure,where,namespace,end,open,abbrev,exact,rw,intro},
  sensitive=true,
  morecomment=[l]{--},
  morecomment=[s]{/-}{-/},
  morestring=[b]",
}
\lstset{
  language=lean,
  basicstyle=\ttfamily\small,
  keywordstyle=\bfseries\color{blue!50!black},
  commentstyle=\itshape\color{green!35!black},
  columns=fullflexible,
  keepspaces=true,
  breaklines=true,
  frame=single,
  framerule=0.3pt,
  rulecolor=\color{black!30},
  xleftmargin=2pt,
  extendedchars=true,
  inputencoding=utf8,
  literate=
    {↔}{{$\leftrightarrow$}}1
    {∃}{{$\exists$}}1
    {∀}{{$\forall$}}1
    {∈}{{$\in$}}1
    {∧}{{$\wedge$}}1
    {≠}{{$\neq$}}1
    {≤}{{$\le$}}1
    {≥}{{$\ge$}}1
    {…}{{$\dots$}}1
    {⟨}{{$\langle$}}1
    {⟩}{{$\rangle$}}1
    {·}{{$\cdot$}}1
    {→}{{$\to$}}1,
}

\title{Proof of Sun's conjecture on coprime adjacent sums\\ (OEIS A228886)}
\author{Claude (AI model by Anthropic)\\ \small run and submitted by Aleksandr Bersenev}
\date{October 8, 2026\\[3pt]\footnotesize Produced autonomously by Claude (Opus 5.5) agents in Claude Code; run and submitted by Aleksandr Bersenev.\\ \footnotesize Not an official Anthropic publication. Not peer reviewed.}
\hypersetup{pdftitle={Proof of Sun's conjecture on coprime adjacent sums (OEIS A228886)},pdfauthor={Claude (AI model by Anthropic); submitted by Aleksandr Bersenev}}

\begin{document}
\maketitle

\begin{abstract}
For $n\ge 1$ let $a(n)$ be the number of permutations $(i_0,i_1,\dots,i_n)$ of $\{0,1,\dots,n\}$ with
$i_0=0$ and $i_n=n$ such that all $n+1$ cyclic adjacent sums
$i_0+i_1,\ i_1+i_2,\ \dots,\ i_{n-1}+i_n,\ i_n+i_0$ are coprime to both $n-1$ and $n+1$
(OEIS A228886). In 2013 Zhi-Wei Sun conjectured that $a(n)>0$ for every $n\ge1$ except $n=2,4$,
and proved this for odd $n$. We give explicit permutations for all even $n\ge 6$, which completes
the proof: $a(n)>0$ if and only if $n\notin\{2,4\}$. We also settle Sun's companion question in which
$n\pm1$ is replaced by $2n\pm1$: a permutation with all cyclic adjacent sums coprime to $2n-1$ and
$2n+1$ exists if and only if $n\neq 2$. In particular the formulation ``for any positive integer $n$''
given in the OEIS entry fails exactly at $n=2$. Both statements, including Sun's odd case and the
impossibility for the exceptional $n$, are formally verified in Lean~4 using the core library only.
\end{abstract}

\tableofcontents

%======================================================================
\section{Introduction}\label{sec:intro}

Sun's 2013 preprint \emph{Some new problems in additive combinatorics}~\cite{Sun1309} contains a
collection of problems asking for permutations or circular arrangements whose adjacent sums or
differences satisfy arithmetic restrictions. Many of them are recorded in the OEIS. One of them is
the following.

\begin{conjecture}[Sun, 2013-09-07; {\cite[Conjecture~3.9]{Sun1309}}, {\cite{A228886}}]\label{conj:sun}
For every positive integer $n\neq 2,4$ there is a permutation $i_0,i_1,\dots,i_n$ of $0,1,\dots,n$ with
$i_0=0$ and $i_n=n$ such that all the $n+1$ adjacent sums
\[
  i_0+i_1,\quad i_1+i_2,\quad\dots,\quad i_{n-1}+i_n,\quad i_n+i_0
\]
are coprime to both $n-1$ and $n+1$.
\end{conjecture}

The number of such permutations is OEIS sequence A228886~\cite{A228886}. Its known values begin
$1,0,1,0,1,8,3,24,78,1164,34,\dots$ for $n=1,2,\dots$, and the entry now lists terms up to $n=26$, all of them positive except
$a(2)=a(4)=0$.
Equivalently, consider the graph on $\{0,\dots,n\}$ in which $i$ and $j$ are joined when $i+j$ is
coprime to $n^2-1$; the conjecture asks for a Hamiltonian cycle of this graph through the edge
$\{0,n\}$.

\paragraph{What was known.}
Sun proved Conjecture~\ref{conj:sun} for all odd $n$ on 2013-09-10, by explicit permutations
\cite[Remark~3.9(ii)]{Sun1309}; see also the comments in~\cite{A228886}. We reproduce his
permutations in \cref{sec:odd}. The even case was left open:
\begin{itemize}
  \item the OEIS entry, edited as recently as July 2026, records only the odd case as proved; its
  attached file lists all permutations for $n\le10$, but no general construction for even $n$;
  \item in a 2020 lecture~\cite{SunTalk2020}, Sun lists the conjecture with the remark ``I have
  proved this for any positive odd integer $n$''.
\end{itemize}

Two related results should not be confused with this one. The analogous problem for permutations of
$1,\dots,n$ whose sums are coprime to $n$ (OEIS A228860~\cite{A228860}) was proved by H.-Q.~Cao and
H.~Pan in 2013, according to that entry. Recent progress on other permutation problems of Sun, such as
the prime circle problem for large $n$ and the Hou--Sun problem on primitive-sum graphs over finite
fields~\cite{She2609}, or the residues of $\sum_k k^2\tau(k)^2$ modulo $2n+1$~\cite{She2606},
does not address A228886.

In Remark~3.9(i) of~\cite{Sun1309}, Sun also guesses that $n\pm1$ in Conjecture~\ref{conj:sun} can be
replaced by $2n\pm1$. The comments of~\cite{A228886} state this companion version as: for any
positive integer $n$ there is such a permutation with all sums coprime to both $2n-1$ and $2n+1$.

\paragraph{Results.}
We prove the even case of Conjecture~\ref{conj:sun} by an explicit construction
(\cref{sec:even}). Together with Sun's odd case and a short argument for $n=2,4$, this gives the
following theorem.

\begin{theorem}\label{thm:main}
For $n\ge1$, a permutation as in Conjecture~\ref{conj:sun} exists if and only if $n\notin\{2,4\}$.
Equivalently, $\mathrm{A228886}(n)>0$ if and only if $n\notin\{2,4\}$.
\end{theorem}

\begin{theorem}\label{thm:companion}
For $n\ge1$, there is a permutation $i_0,\dots,i_n$ of $0,\dots,n$ with $i_0=0$, $i_n=n$ and all
$n+1$ cyclic adjacent sums coprime to both $2n-1$ and $2n+1$ if and only if $n\neq2$.
\end{theorem}

For $n=2$ the only candidate is $(0,1,2)$, and $1+2=3=2n-1$. So the OEIS wording of the companion
conjecture (``for any positive integer $n$'') is false at $n=2$ and true for every other $n$. The
wording in~\cite{Sun1309} inherits the exceptions $n\ne2,4$ of Conjecture~\ref{conj:sun} and is
therefore true.

Both theorems are formally verified in Lean~4 (\cref{sec:lean}). The verification covers the odd
case, the even case, the companion problem and the impossibility at the exceptional values.

\paragraph{Novelty.}
To our knowledge, as of October 8, 2026, neither the even case of Conjecture~\ref{conj:sun} nor the
companion conjecture had been proved. Our search covered:
\begin{itemize}
  \item the OEIS entry and related entries;
  \item the latest arXiv version (v9, March 2020) of~\cite{Sun1309}, in which Conjecture~3.9 and
  Remark~3.9 still record only the odd case, and Sun's 2020 slides~\cite{SunTalk2020};
  \item web and arXiv searches with several phrasings, in English and Chinese;
  \item the OEIS problem list of the \emph{formal-conjectures} project~\cite{FormalConj}, which
  does not contain A228886, and the lists of OEIS conjectures reported in 2026 as proved by AI
  agents~\cite{AI2605,AI2608}, which do not mention it either.
\end{itemize}
We could \emph{not} consult Sun's 2021 Chinese-language book of conjectures~\cite{SunBook}. If the
even case was settled there or anywhere else, our contribution reduces to the companion result and
the formal verification. The constructions are short and elementary. The even case seems to have
stayed open for lack of attention rather than difficulty.

%======================================================================
\section{Definitions and notation}\label{sec:defs}

Throughout, $n\ge1$ is an integer, and $\gcd(s,0)=s$, so that being coprime to $0$ means being $1$.

\begin{definition}
Let $A,B\ge0$ be integers. An \emph{$(A,B)$-cycle of order $n$} is a sequence
$p=(p(0),p(1),\dots,p(n))$ such that:
\begin{enumerate}
  \item $p$ is a permutation of $\{0,1,\dots,n\}$;
  \item $p(0)=0$ and $p(n)=n$;
  \item every cyclic adjacent sum
  \[
    s_i:=p(i)+p(i+1)\quad(0\le i\le n-1),\qquad s_n:=p(n)+p(0)=n,
  \]
  is coprime to both $A$ and $B$.
\end{enumerate}
\end{definition}

In this language, Theorem~\ref{thm:main} concerns $(n-1,n+1)$-cycles and Theorem~\ref{thm:companion} concerns
$(2n-1,2n+1)$-cycles. We call an integer $s$ \emph{admissible} when it is coprime to both moduli.
Permutations are written as their value sequences $(p(0),\dots,p(n))$.

The following triviality is used repeatedly. If $d$ divides $a$ and $b$, then $d$ divides every
integer combination $xa+yb$. In particular $\gcd(a,a+k)$ divides $k$.

%======================================================================
\section{Admissible sums}\label{sec:admissible}

\begin{lemma}[moduli $n\pm1$]\label{lem:adm1}
Let $n\ge1$. The following sums are coprime to both $n-1$ and $n+1$:
\begin{enumerate}
  \item $1$ and $n$, always;
  \item $n+3$, if $n$ is even; and $n-3$, if $n$ is even and $n\ge4$;
  \item $n+2$, if $n\not\equiv1\pmod 3$;
  \item $n-2$ (for $n\ge3$) and $2n-1$, if $n\not\equiv2\pmod 3$.
\end{enumerate}
\end{lemma}

\begin{proof}
\emph{(1)} For $1$ this is clear. For $n$, note $\gcd(n,n\pm1)\mid 1$.

\emph{(2)} If $n$ is even, then $n\pm1$ and $n\pm3$ are odd. Now $\gcd(n+3,n+1)\mid 2$ and
$\gcd(n+3,n-1)\mid4$; since all four numbers are odd, both gcds are $1$. Similarly
$\gcd(n-3,n-1)\mid2$ and $\gcd(n-3,n+1)\mid 4$.

\emph{(3)} We have $\gcd(n+2,n+1)=1$ and $\gcd(n+2,n-1)\mid 3$. If $3\mid n+2$ then $n\equiv1$, so for
$n\not\equiv1\pmod3$ the second gcd is $1$.

\emph{(4)} We have $\gcd(n-2,n-1)=1$ and $\gcd(n-2,n+1)\mid3$, and $3\nmid n-2$ when
$n\not\equiv2$. Also $2n-1=2(n-1)+1$ gives $\gcd(2n-1,n-1)=1$. Finally $2(n+1)-(2n-1)=3$ gives
$\gcd(2n-1,n+1)\mid3$, and $3\mid 2n-1$ iff $n\equiv2\pmod 3$.
\end{proof}

\begin{lemma}[moduli $2n\pm1$]\label{lem:adm2}
Let $n\ge1$. The following sums are coprime to both $2n-1$ and $2n+1$:
\begin{enumerate}
  \item $n$ and $4$, always; and $2n-3$, if $n\ge2$;
  \item $n/2$, if $n$ is even;
  \item $n+1$, and $(n+1)/2$ for odd $n$, if $n\not\equiv2\pmod3$;
  \item $n-1$ (for $n\ge2$), and $(n-1)/2$ for odd $n$, if $n\not\equiv1\pmod 3$.
\end{enumerate}
\end{lemma}

\begin{proof}
Both moduli are odd.

\emph{(1)} $\gcd(n,2n\pm1)\mid 1$. The number $4$ is coprime to odd numbers.
$\gcd(2n-3,2n-1)\mid2$ and $\gcd(2n-3,2n+1)\mid 4$, and $2n-3$ is odd.

\emph{(2)} With $m=n/2$: $\gcd(m,4m\pm1)\mid1$.

\emph{(3)} $\gcd(n+1,2n+1)\mid 2(n+1)-(2n+1)=1$. Also $\gcd(n+1,2n-1)\mid 3$, and $3\mid n+1$ iff
$n\equiv2$. For odd $n$ put $x=(n+1)/2$, so $2n-1=4x-3$ and $2n+1=4x-1$. Then $\gcd(x,4x-1)=1$ and
$\gcd(x,4x-3)\mid3$. Moreover $3\mid x$ iff $n\equiv5\pmod 6$, which for odd $n$ means
$n\equiv2\pmod3$.

\emph{(4)} $\gcd(n-1,2n-1)\mid (2n-1)-2(n-1)=1$. Also $\gcd(n-1,2n+1)\mid3$, and $3\mid n-1$ iff
$n\equiv1$. For odd $n$ put $y=(n-1)/2$, so $2n-1=4y+1$ and $2n+1=4y+3$. Then $\gcd(y,4y+1)=1$ and
$\gcd(y,4y+3)\mid3$. Moreover $3\mid y$ iff $n\equiv1\pmod6$, which for odd $n$ means
$n\equiv1\pmod 3$.
\end{proof}

%======================================================================
\section{Proof of \texorpdfstring{Theorem~\ref{thm:main}}{Theorem 1.2}}\label{sec:proof1}

\subsection{The exceptional values \texorpdfstring{$n=2,4$}{n = 2, 4}}

\begin{proposition}\label{prop:24}
There is no $(n-1,n+1)$-cycle of order $2$ or of order $4$.
\end{proposition}
\begin{proof}
For $n=2$ the only permutation with $p(0)=0$ and $p(2)=2$ is $(0,1,2)$. Its sum $1+2=3$ is not
coprime to $n+1=3$.

For $n=4$ write the permutation as $(0,a,b,c,4)$ with $\{a,b,c\}=\{1,2,3\}$; the moduli are $3$
and $5$. The sum $c+4$ must avoid $5$ and $6$, so $c=3$. Then $\{a,b\}=\{1,2\}$, and $a+b=3$ is
not coprime to $3$.
\end{proof}

\subsection{Odd \texorpdfstring{$n$}{n} (Sun's permutations)}\label{sec:odd}

The following two families are due to Sun~\cite[Remark~3.9(ii)]{Sun1309}. We record them with
their sums.

\begin{construction}[Sun]\label{con:oddA}
Let $n$ be odd with $n\not\equiv2\pmod3$. Define
\[
  p(i)=\begin{cases} i & i \text{ even},\ i\le n-1,\\ n-1-i & i\text{ odd},\\ n & i=n.\end{cases}
\]
This gives $(0,\,n-2,\,2,\,n-4,\,4,\,\dots,\,1,\,n-1,\,n)$. The sums are:
\begin{itemize}
  \item $i+(n-2-i)=n-2$ for even $i<n-1$;
  \item $(n-1-i)+(i+1)=n$ for odd $i$;
  \item $s_{n-1}=(n-1)+n=2n-1$ and $s_n=n$.
\end{itemize}
\end{construction}

\begin{construction}[Sun]\label{con:oddB}
Let $n\ge3$ be odd with $n\not\equiv1\pmod3$. Define $p(0)=0$, $p(n)=n$, $p(i)=i$ for odd $i$, and
$p(i)=n+1-i$ for even $2\le i\le n-1$. This gives $(0,\,1,\,n-1,\,3,\,n-3,\,\dots,\,n-2,\,2,\,n)$.
The sums are:
\begin{itemize}
  \item $s_0=1$;
  \item $n$ for odd $i$;
  \item $n+2$ for even $2\le i\le n-1$;
  \item $s_n=n$.
\end{itemize}
\end{construction}

Both maps are permutations. In Construction~\ref{con:oddA} even positions carry even values and
odd positions carry odd values; in Construction~\ref{con:oddB} even positions $i\ge2$ carry the even
values $n+1-i$. By Lemma~\ref{lem:adm1}(1),(3),(4) all sums are admissible. This proves Theorem~\ref{thm:main}
for odd $n$, including $n=1$, where the permutation is $(0,1)$ with sums $1,1$.

\subsection{Even \texorpdfstring{$n$}{n} (new)}\label{sec:even}

Let $n=2m$ be even. Sun's zig-zag $i\mapsto i$ ($i$ odd), $i\mapsto n+1-i$ ($i$ even) has adjacent sums
alternating between $n$ and $n+2$. For even $n$ it cannot end at $p(n)=n$, because it sends the even
position $n$ to the value $1$. Our permutations therefore switch, in the middle, to the opposite
phase $i\mapsto i$ ($i$ even), $i\mapsto n+1-i$ ($i$ odd), using a short gadget whose sums are
$n\pm3$, which are always admissible for even $n$ (Lemma~\ref{lem:adm1}(2)).

\begin{construction}[$n\equiv0\pmod4$, $n\ge8$]\label{con:evenA}
Let $m=n/2$, which is even. Define
\[
p(i)=\begin{cases}
0 & i=0,\\
i & 1\le i\le m-3,\ i\text{ odd},\\
n+1-i & 1\le i\le m-3,\ i \text{ even},\\
m,\ m+3,\ m-1,\ m+1 & i=m-2,\ m-1,\ m,\ m+1 \text{ respectively},\\
i & m+2\le i\le n,\ i\text{ even},\\
n+1-i & m+2\le i\le n,\ i\text{ odd}.
\end{cases}
\]
That is, $p$ is the sequence
\[
\begin{gathered}0,\ 1,\ n-1,\ 3,\ n-3,\ \dots,\ m+5,\ m-3,\ \big|\ m,\ m+3,\ m-1,\ m+1\ \big|\\ m+2,\ m-2,\ m+4,\ m-4,\ \dots,\ n-2,\ 2,\ n.\end{gathered}
\]
\end{construction}

\begin{construction}[$n\equiv2\pmod4$, $n\ge10$]\label{con:evenB}
Let $m=n/2$, which is odd. Define
\[
p(i)=\begin{cases}
0 & i=0,\\
i & 1\le i\le m-1,\ i\text{ odd},\\
n+1-i & 1\le i\le m-1,\ i \text{ even},\\
m+1,\ m-1,\ m+3,\ m,\ m-3 & i=m,\ m+1,\ m+2,\ m+3,\ m+4 \text{ respectively},\\
i & m+5\le i\le n,\ i\text{ even},\\
n+1-i & m+5\le i\le n,\ i\text{ odd}.
\end{cases}
\]
That is, $p$ is the sequence
\[
0,\ 1,\ n-1,\ 3,\ n-3,\ \dots,\ m-2,\ m+2,\ \big|\ m+1,\ m-1,\ m+3,\ m,\ m-3\ \big|\ m+5,\ m-5,\ \dots,\ n-2,\ 2,\ n.
\]
\end{construction}

\begin{lemma}\label{lem:evenAB}
Constructions~\ref{con:evenA} and~\ref{con:evenB} define permutations of $\{0,\dots,n\}$ with $p(0)=0$ and $p(n)=n$. All
their cyclic adjacent sums lie in $\{1,\ n,\ n+2,\ n+3,\ n-3\}$. Precisely:
\begin{itemize}
  \item Construction~\ref{con:evenA}:
  \begin{itemize}
    \item $s_0=1$;
    \item $s_i=n$ ($i$ odd) and $s_i=n+2$ ($i$ even) for $1\le i\le m-4$;
    \item $(s_{m-3},s_{m-2},s_{m-1},s_m,s_{m+1})=(n-3,\,n+3,\,n+2,\,n,\,n+3)$;
    \item $s_i=n$ ($i$ even) and $s_i=n+2$ ($i$ odd) for $m+2\le i\le n-1$;
    \item $s_n=n$.
  \end{itemize}
  \item Construction~\ref{con:evenB}:
  \begin{itemize}
    \item $s_0=1$;
    \item $s_i=n$ ($i$ odd) and $s_i=n+2$ ($i$ even) for $1\le i\le m-2$;
    \item $(s_{m-1},\dots,s_{m+4})=(n+3,\,n,\,n+2,\,n+3,\,n-3,\,n+2)$;
    \item $s_i=n$ ($i$ even) and $s_i=n+2$ ($i$ odd) for $m+5\le i\le n-1$;
    \item $s_n=n$.
  \end{itemize}
\end{itemize}
\end{lemma}

\begin{proof}
\emph{Permutation, Construction~\ref{con:evenA}.} Here $m$ is even. We list the values taken by each block.
\begin{itemize}
  \item First block, $1\le i\le m-3$. The odd positions take the odd values $1,3,\dots,m-3$. The even
  positions $2\le i\le m-4$ take the values $n+1-i$, which are the odd values $m+5,\dots,n-1$.
  \item Gadget. It takes $m$ and the odd values $m-1,m+1,m+3$.
  \item Last block, $m+2\le i\le n$. The even positions take the even values $m+2,\dots,n$. The odd
  positions $m+3\le i\le n-1$ take the values $n+1-i$, which are the even values $2,\dots,m-2$.
\end{itemize}
Together with $p(0)=0$ every odd value in $[1,n-1]$ and every even value in $[0,n]$ occurs exactly
once.

\emph{Permutation, Construction~\ref{con:evenB}.} Here $m$ is odd and $m\ge5$.
\begin{itemize}
  \item First block, $1\le i\le m-1$. It takes the odd values $1,\dots,m-2$ and, at even positions,
  the odd values $n+1-i\in\{m+2,\dots,n-1\}$.
  \item Gadget. It takes the odd value $m$ and the even values $m-3,m-1,m+1,m+3$.
  \item Last block, $m+5\le i\le n$. Here $m+5$ is even. It takes the even values $m+5,\dots,n$
  and, at odd positions, the even values $n+1-i\in\{2,\dots,m-5\}$.
\end{itemize}
Again every value in $[0,n]$ occurs exactly once.

\emph{Sums.} Inside a block of the first phase, an odd position $i$ gives
$i+\big(n+1-(i+1)\big)=n$ and an even position gives $(n+1-i)+(i+1)=n+2$. The same holds with the
parities exchanged in a block of the second phase.

For Construction~\ref{con:evenA}, the transitions around the gadget are
\[
\begin{gathered}(m-3)+m=n-3,\quad m+(m+3)=n+3,\quad (m+3)+(m-1)=n+2,\\ (m-1)+(m+1)=n,\quad (m+1)+(m+2)=n+3.\end{gathered}
\]
For Construction~\ref{con:evenB} they are
\[
\begin{gathered}
(m+2)+(m+1)=n+3,\quad (m+1)+(m-1)=n,\quad (m-1)+(m+3)=n+2,\\
(m+3)+m=n+3,\quad m+(m-3)=n-3,\quad (m-3)+(m+5)=n+2.
\end{gathered}
\]
Here $p(m-1)=n+1-(m-1)=m+2$, because $m-1$ is even. Finally $s_0=0+1=1$ and $s_n=n+0=n$.
\end{proof}

\begin{lemma}[mirror]\label{lem:mirror}
If $p$ is a permutation of $\{0,\dots,n\}$ with $p(0)=0$ and $p(n)=n$, then so is
$q(i):=n-p(n-i)$. Its adjacent sums are
\[
q(i)+q(i+1)=2n-s_{n-1-i}\quad(0\le i\le n-1),\qquad q(n)+q(0)=n .
\]
\end{lemma}
\begin{proof}
The map $i\mapsto n-i$ reverses the order and the map $v\mapsto n-v$ reflects the values; both are
involutions. We have $q(0)=n-p(n)=0$ and $q(n)=n-p(0)=n$. The sum formula is immediate:
\[
q(i)+q(i+1)=2n-\big(p(n-i)+p(n-i-1)\big)=2n-s_{n-1-i}. \qedhere
\]
\end{proof}

\begin{proposition}\label{prop:even}
For every even $n\ge6$ there is an $(n-1,n+1)$-cycle of order $n$.
\end{proposition}
\begin{proof}
For $n=6$ take $(0,1,2,4,5,3,6)$. Its cyclic sums are $1,3,6,9,8,9,6$, all coprime to $5$ and $7$.

Let $n\ge8$ be even, and let $p$ be Construction~\ref{con:evenA} if $4\mid n$ and Construction~\ref{con:evenB} otherwise.
\begin{itemize}
  \item If $n\not\equiv1\pmod3$, all sums of $p$ lie in $\{1,n,n+2,n\pm3\}$ by Lemma~\ref{lem:evenAB}.
  These are admissible by Lemma~\ref{lem:adm1}(1)--(3).
  \item If $n\equiv1\pmod3$, use the mirror $q$ of $p$ (Lemma~\ref{lem:mirror}). Its sums lie in
  $\{2n-1,\,n,\,n-2,\,n\mp3\}$. These are admissible by Lemma~\ref{lem:adm1}(1),(2),(4), because
  $n\not\equiv2\pmod 3$. \qedhere
\end{itemize}
\end{proof}

\begin{proof}[Proof of Theorem~\ref{thm:main}]
Combine Proposition~\ref{prop:24}, Sun's odd constructions (\cref{sec:odd}) and Proposition~\ref{prop:even}.
\end{proof}

\begin{example}\label{ex:even}
The permutations produced by the proof for small even $n$ are listed below. For $n=10$ and $n=16$
the mirror is used, because $n\equiv1\pmod3$.
\begin{center}\footnotesize
\begin{tabular}{@{}rll@{}}
\toprule
$n$ & permutation $(p(0),\dots,p(n))$ & cyclic sums $(s_0,\dots,s_n)$\\
\midrule
6 & $0,1,2,4,5,3,6$ & $1,3,6,9,8,9,6$\\
8 & $0,1,4,7,3,5,6,2,8$ & $1,5,11,10,8,11,8,10,8$\\
10 & $0,8,5,2,6,4,3,7,1,9,10$ & $8,13,7,8,10,7,10,8,10,19,10$\\
12 & $0,1,11,3,6,9,5,7,8,4,10,2,12$ & $1,12,14,9,15,14,12,15,12,14,12,14,12$\\
14 & $0,1,13,3,11,5,9,8,6,10,7,4,12,2,14$ & $1,14,16,14,16,14,17,14,16,17,11,16,14,16,14$\\
16 & $0,14,2,12,4,10,6,7,9,5,8,11,3,13,1,15,16$ & $14,16,14,16,14,16,13,16,14,13,19,14,16,14,16,31,16$\\
\bottomrule
\end{tabular}
\end{center}
\end{example}

\begin{remark}[how the construction was found]
Put $t=v-m\in[-m,m]$. A sum $v+w=n+c$ becomes $t+t'=c$. The always-admissible sums $n$ and $n\pm3$
correspond to the reflections $t\mapsto -t$ and $t\mapsto\pm3-t$. Adding $n+2$, which corresponds to
$t\mapsto 2-t$, makes the two zig-zags $0,2,-2,4,-4,\dots$ and $1,-1,3,-3,\dots$ available. They are
glued near $t=0$ by the gadget $2,1,-1,3,0,-3,5$, and closed near $t=\pm m$ by the sum $v+w=1$. The
mirror $t\mapsto-t$ exchanges $n+2$ with $n-2$, which covers the residue class $n\equiv1\pmod 3$. An exhaustive search shows that for $n=8$ there is no admissible cycle using only the sums $n$,
$n\pm3$ and powers of two, so these universally admissible sums alone do not suffice.
\end{remark}

%======================================================================
\section{Proof of \texorpdfstring{Theorem~\ref{thm:companion}}{Theorem 1.3}}\label{sec:companion}

Now the moduli are $2n-1$ and $2n+1$.

\begin{construction}[$n\not\equiv2\pmod3$, $n$ even]\label{con:C1}
Let $m=n/2$, $p(0)=0$ and $p(n)=n$. For $1\le i\le n-1$ put
$p(i)=m-\tfrac{i-1}{2}$ for odd $i$ and $p(i)=m+\tfrac i2$ for even $i$. This gives
\[
(0,\ m,\ m+1,\ m-1,\ m+2,\ m-2,\ \dots,\ n-1,\ 1,\ n).
\]
The odd positions take $\{1,\dots,m\}$ and the even positions $2\le i\le n-2$ take
$\{m+1,\dots,n-1\}$. The sums are:
\begin{itemize}
  \item $s_0=m$;
  \item $s_i=n+1$ for odd $i\le n-3$ and $s_i=n$ for even $2\le i\le n-2$;
  \item $s_{n-1}=1+n$ and $s_n=n$.
\end{itemize}
\end{construction}

\begin{construction}[$n\not\equiv2\pmod3$, $n$ odd]\label{con:C2}
Let $m=(n-1)/2$, $p(0)=0$ and $p(n)=n$. For $1\le i\le n-1$ put
$p(i)=m+1+\tfrac{i-1}2$ for odd $i$ and $p(i)=m+1-\tfrac i2$ for even $i$. This gives
\[
(0,\ m+1,\ m,\ m+2,\ m-1,\ \dots,\ n-1,\ 1,\ n).
\]
The odd positions take $\{m+1,\dots,n-1\}$ and the even positions take $\{1,\dots,m\}$. The sums are:
\begin{itemize}
  \item $s_0=(n+1)/2$;
  \item $s_i=n$ for odd $i$ and $s_i=n+1$ for even $2\le i\le n-3$;
  \item $s_{n-1}=n+1$ and $s_n=n$.
\end{itemize}
For $n=1$ this is $(0,1)$.
\end{construction}

\begin{construction}[$n\equiv2\pmod3$, $n\ge8$]\label{con:C3}
Let $m=\lfloor n/2\rfloor$ and $p(0)=0$. For $1\le i\le n-6$ put:
\begin{itemize}
  \item if $n$ is even: $p(i)=m+\tfrac{i-1}2$ for odd $i$, and $p(i)=m-\tfrac i2$ for even $i$;
  \item if $n$ is odd: $p(i)=m-\tfrac{i-1}2$ for odd $i$, and $p(i)=m+\tfrac i2$ for even $i$.
\end{itemize}
Finally put $(p(n-5),p(n-4),\dots,p(n))=(1,\,n-1,\,n-2,\,2,\,n-3,\,n)$. So $p$ is
\[
0,\ \underbrace{m,\ m\mp1,\ m\pm1,\ \dots,\ 4,\ n-4,\ 3}_{\text{positions }1,\dots,n-6},\ 1,\ n-1,\ n-2,\ 2,\ n-3,\ n,
\]
with the upper signs for even $n$. The zig-zag block takes exactly the values $\{3,\dots,n-4\}$. The sums are:
\begin{itemize}
  \item $s_0=m$;
  \item along the zig-zag, $n-1$ and $n$ alternately;
  \item $s_{n-6}=3+1=4$;
  \item then $n,\ 2n-3,\ n,\ n-1,\ 2n-3$;
  \item $s_n=n$.
\end{itemize}
\end{construction}

\begin{proof}[Verification of Construction~\ref{con:C3}]
\emph{$n$ even.} The odd positions $1,3,\dots,n-7$ take the values $m,m+1,\dots,m+\frac{n-8}2=n-4$.
The even positions $2,\dots,n-6$ take the values $m-1,\dots,m-\frac{n-6}2=3$. An odd position
$i=2k+1$ gives $(m+k)+(m-k-1)=n-1$, and an even position $i=2k$ gives $(m-k)+(m+k)=n$.

\emph{$n$ odd.} Here $n=2m+1$. The odd positions take $m,m-1,\dots,3$ and the even positions take
$m+1,\dots,n-4$. The zig-zag sums are $(m-k)+(m+k+1)=n$ and $(m+k)+(m-k)=n-1$.

In both cases the last zig-zag value is $p(n-6)=3$, and the tail gives the sums listed.
\end{proof}

\begin{proof}[Proof of Theorem~\ref{thm:companion}]
For $n=2$ the only candidate $(0,1,2)$ has the sum $3=2n-1$, so no such permutation exists.

For $n=5$ take $(0,1,4,3,2,5)$. Its sums $1,5,7,5,7,5$ are coprime to $9$ and $11$.

For all other $n$ we use the constructions above.
\begin{itemize}
  \item If $n\not\equiv2\pmod 3$, use Construction~\ref{con:C1} or Construction~\ref{con:C2}. All sums lie in
  $\{n/2,\,n,\,n+1\}$ or $\{(n+1)/2,\,n,\,n+1\}$, which are admissible by Lemma~\ref{lem:adm2}(1)--(3).
  \item If $n\equiv2\pmod 3$ and $n\ge8$, use Construction~\ref{con:C3}. All sums lie in
  $\{m,\,n-1,\,n,\,4,\,2n-3\}$, which are admissible by Lemma~\ref{lem:adm2}(1),(2),(4). Here, for odd $n$,
  $m=(n-1)/2$ and $n\not\equiv1\pmod 3$. \qedhere
\end{itemize}
\end{proof}

\begin{example}
For $n=4$, which the arXiv formulation of the companion conjecture excludes,
Construction~\ref{con:C1} gives $(0,2,3,1,4)$ with sums $2,5,4,5,4$
(moduli $7,9$). Further examples:
\begin{itemize}
  \item $n=8$: $(0,4,3,1,7,6,2,5,8)$, with sums $4,7,4,8,13,8,7,13,8$ (moduli $15,17$);
  \item $n=11$: $(0,5,6,4,7,3,1,10,9,2,8,11)$, with sums $5,11,10,11,10,4,11,19,11,10,19,11$
  (moduli $21,23$).
\end{itemize}
\end{example}

%======================================================================
\section{Formal verification}\label{sec:lean}

All results above are verified in a single self-contained Lean~4~\cite{Lean4} file, \texttt{lean/Main.lean}
(about 830 lines). It uses Lean version 4.34.1 with the core library only: no Mathlib and no
external dependencies. It contains no \texttt{sorry}, \texttt{axiom}, \texttt{admit} or
\texttt{native\_decide}.

\paragraph{Statement.}
The two main theorems are stated as follows. Their definitions appear directly above them in the
file and are reproduced here verbatim.
\lstinputlisting{sections/leandefs.lean}
\lstinputlisting{sections/leanthm.lean}

\paragraph{Axioms and resources.}
\texttt{\#print axioms} reports:
\begin{lstlisting}[language={}]
'A228886.sun_conjecture' depends on axioms: [propext, Classical.choice, Quot.sound]
'A228886.sun_companion' depends on axioms: [propext, Classical.choice, Quot.sound]
\end{lstlisting}
These are the standard axioms of Lean's foundations.
\begin{itemize}
  \item The check takes about 480\,s of wall-clock time (about 352\,s of CPU time) with a peak
  memory of about 1.6\,GB, on a shared machine.
  \item The file sets \texttt{set\_option Elab.async false}, because asynchronous elaboration failed
  to create threads under that machine's resource limits. The option has no logical effect.
\end{itemize}

\paragraph{Statement audit.}
\begin{itemize}
  \item \texttt{List.range (n + 1)} is $[0,1,\dots,n]$. So \texttt{l.Perm (List.range (n + 1))}
  says that $l$ lists $0,\dots,n$, each exactly once.
  \item \texttt{l.head? = some 0} and \texttt{l.getLast? = some n} say $i_0=0$ and $i_n=n$.
  \item For a list of length $n+1$, \texttt{cyclicSums l} is the list of
  $l[i]+l[(i+1)\bmod(n+1)]$ for $i=0,\dots,n$, that is $i_0+i_1,\dots,i_{n-1}+i_n,i_n+i_0$. Both
  indices are always smaller than the length, so the default value $0$ of \texttt{getD} is never
  used.
  \item \texttt{Nat.gcd s A = 1} is coprimality. Under the hypothesis $1\le n$ the natural-number
  subtractions $n-1$ and $2n-1$ are exact. For $n=1$, $\gcd(s,0)=s$, matching the OEIS value
  $a(1)=1$.
  \item The theorems are equivalences for every $n\ge1$, which is the offset of A228886. They include
  the non-existence statements at $n=2,4$ (respectively $n=2$).
  \item The file also contains \texttt{decide}-checked examples. They show that a correct
  permutation is accepted, and that a list violating the gcd condition, a rotated cycle (same cyclic
  sums but $i_0\neq0$), and a non-permutation are all rejected.
\end{itemize}

\paragraph{Proof structure in Lean.}
\begin{itemize}
  \item An equivalent function form \texttt{SunFun n A B p} is used, in which
  $p:\N\to\N$ is injective on $\{0,\dots,n\}$ with values in $\{0,\dots,n\}$. The lemma
  \texttt{exists\_list\_iff} transfers between the list and function forms, using a pigeonhole
  lemma \texttt{surj\_of\_inj} proved by induction.
  \item Each construction of \cref{sec:proof1,sec:companion} is a Lean function defined by a cascade
  of \texttt{if}s: \texttt{oddA}, \texttt{oddB}, \texttt{evenA}, \texttt{evenB}, the mirror,
  \texttt{compA}, \texttt{compB}, \texttt{compC}, and the explicit cases $n=5,6$.
  \item Bounds, injectivity and the list of adjacent sums are proved by splitting all branches and
  calling \texttt{omega}; injectivity is checked over all pairs of branches.
  \item The admissibility lemmas of \cref{sec:admissible} are reduced, by Euclid-type rewriting, to
  $\gcd(r,k)=1$ with $k\le 4$, which is settled by \texttt{decide}.
  \item The impossibility at $n=2,4$ is a short case analysis.
\end{itemize}

\paragraph{Independent numerical check.}
The script \texttt{data/verify\_constructions.py} transcribes the Lean definitions into Python. It
confirms every construction for $n\le 3000$ and recomputes $\mathrm{A228886}(1..11)$ by brute force, in
agreement with the OEIS data. This check is not part of the proof.

%======================================================================
\section{Limitations and open questions}\label{sec:limits}

\begin{itemize}
\item \textbf{Novelty.} The claim of novelty rests on a literature and web search, as described in
\cref{sec:intro}. In particular we could not check Sun's 2021 book~\cite{SunBook}.
\item \textbf{Depth.} The proofs are elementary. They give existence only and say nothing about the
size of $\mathrm{A228886}(n)$, which grows irregularly (for instance $a(11)=34$ while
$a(12)=4021156$). Its dependence on the factorisation of $n^2-1$ is not addressed.
\item \textbf{Squares.} Sun's variant A229130~\cite{A229130} uses $i_k^2+i_{k+1}$ in place of
$i_k+i_{k+1}$. It conjectures existence for all $n\notin\{2,5,11\}$, and more generally that the
exceptional sets $E(k)$ for $i_k^{\,k}+i_{k+1}$ are finite. It remains open as far as we know.
The linear zig-zag method used here does not apply directly.
\item \textbf{Other problems in~\cite{Sun1309}.} Other circular-permutation problems there involve
primes (twin primes, Sophie Germain primes, prime circles) or primitive roots. They are of a
different nature, and most remain open.
\end{itemize}

\begin{thebibliography}{99}
\bibitem{A228886}
OEIS Foundation Inc., \emph{The On-Line Encyclopedia of Integer Sequences}, sequence A228886
(Z.-W. Sun, Sep 07 2013; terms to $a(26)$ by M. Alekseyev, Jul 07 2026).
\url{https://oeis.org/A228886}.

\bibitem{A228860}
OEIS Foundation Inc., sequence A228860 (Z.-W. Sun, Sep 05 2013).
\url{https://oeis.org/A228860}.

\bibitem{A229130}
OEIS Foundation Inc., sequence A229130 (Z.-W. Sun, Sep 15 2013).
\url{https://oeis.org/A229130}.

\bibitem{Sun1309}
Z.-W. Sun, \emph{Some new problems in additive combinatorics}, Nanjing Univ. J. Math. Biquarterly
36 (2019), 134--155; arXiv:1309.1679 [math.NT], v9 (March 2, 2020) consulted.
\url{https://arxiv.org/abs/1309.1679}.

\bibitem{SunTalk2020}
Z.-W. Sun, \emph{Introduction to Combinatorial Number Theory (VI): Permutations, Determinants and
Permanents}, on-line talk slides, June 26, 2020.
\url{http://maths.nju.edu.cn/~zwsun/CNT-talk6.pdf}.

\bibitem{SunBook}
Z.-W. Sun, \emph{New Conjectures in Number Theory and Combinatorics} (in Chinese), Harbin Institute
of Technology Press, 2021. (Not consulted for this work.)

\bibitem{She2606}
Y.-F. She and X.-Q. Luo, \emph{Proof of a conjecture on permutations}, arXiv:2606.18311, 2026.

\bibitem{She2609}
Y.-F. She, \emph{Hamiltonicity in graphs defined by primes and primitive elements},
arXiv:2609.19114, 2026.

\bibitem{FormalConj}
Google DeepMind, \emph{formal-conjectures} (OEIS folder),
\url{https://github.com/google-deepmind/formal-conjectures}, accessed October 8, 2026.

\bibitem{AI2605}
G. Tsoukalas, A. Kovsharov, S. Shirobokov, A. Surina, M. Firsching, G. B\'erczi, F. J. R. Ruiz,
A. Suggala, A. Z. Wagner, E. Wieser, L. Yu, A. Huang, et al., \emph{Advancing mathematics research
with AI-driven formal proof search}, arXiv:2605.22763, 2026.

\bibitem{AI2608}
T. Adamczewski, \emph{OEIS Open: How many conjectures can language models turn into theorems?},
arXiv:2608.11941, 2026.

\bibitem{Lean4}
L.~de Moura and S.~Ullrich, \emph{The Lean 4 theorem prover and programming language}, in
Automated Deduction -- CADE 28, Lecture Notes in Computer Science 12699, Springer, 2021, 625--635.
\end{thebibliography}

\end{document}
